The $1/x^2$ potential with the naive boundary condition: $\psi(x) \to 0$ for $x \to 0$, poses clear challenges. It leads to an infinite and uncountable number of bound states for negative energies, unbounded in absolute value.

\vspace{4mm}
\noindent
Alternatively, one could propose that some potentials in Quantum Mechanics are simply forbidden or nonexistent in the natural world, but I do not agree. There is experimental evidence indicating the existence of physical systems represented, at least approximately, by the $1/x^2$ potential or its three-dimensional analog. A prime example is the interaction of an electron with a stationary electric dipole, such as an electron interacting with a polar molecule. In this case, the potential is $-e p \cos(\theta)/r^2$, and the radial part of the Schrödinger equation aligns mathematically with our discussion. Furthermore, the critical parameter $\alpha = 1/4$ holds significance in the study of this system, as it is associated with the minimum value for the dipole moment or the minimum moment required to bind a charged particle to an extended electric dipole.

\vspace{4mm}
\noindent
In conclusion, the discussion presented in this essay might appear somewhat reminiscent of non-ordinary quantum mechanics. However, this resemblance is likely due to the fact that introductory quantum mechanics courses often teach that certain mathematical intricacies are the domain of mathematicians. This is not always the case, which is why I believe a rigorous mathematical approach to physics should be as fundamental as the teaching of intuition.